\documentclass[11pt,a4paper]{article}
\usepackage[T1]{fontenc}
\usepackage{lmodern}
\usepackage[margin=24mm]{geometry}
\usepackage{amsmath,amssymb,booktabs,microtype}
\usepackage[hidelinks,hypertexnames=false]{hyperref}
\hypersetup{pdftitle={Hydrogen escape: mathematical model definitions},
  pdfauthor={exojax-escape}}
% Omit volatile metadata so rebuilding with the same TeX toolchain is stable.
\pdfinfoomitdate=1
\pdftrailerid{}
\pdfsuppressptexinfo=15
\newcommand{\dd}{\mathrm{d}}
\newcommand{\nH}{n_{\mathrm H}}
\newcommand{\kb}{k_{\mathrm B}}
\newcommand{\mpart}{m_{\mathrm p}}
\newcommand{\wn}{\tilde\nu}
\newcommand{\Rp}{R_{\mathrm p}}
\newcommand{\status}[1]{\par\noindent\textbf{Status: #1.}\ }
\title{Hydrogen escape: mathematical model definitions}
\author{exojax-escape}
\date{Definition revision 1 --- 20 November 2026}

\begin{document}
\maketitle

\section{Scope, status and conventions}\label{sec:scope}
This is the mathematical specification accompanying the code and its MCEG
records. Implemented equations below describe source revision
\href{https://github.com/HajimeKawahara/exojax-escape/tree/421f6a532ced0d0c5614124c4fea902070f93174}{\texttt{421f6a5}}.
The proposed equations specify the next NLTE development stage; they are not
an implemented solver or a reproduction of the atmosphere in Zhang et
al.~\cite{zhang}. Model choices still needed to close that problem are listed
in Section~\ref{sec:choices}. Empirical comparison Models 1 and 2 remain in the
inventory but are outside this document's initial equation scope.

\begin{center}
\begin{tabular}{@{}ll@{}}
\toprule
Component & Status in the audited implementation \\
\midrule
Pure-H continuity; isolated two-level ratio & Implemented \\
Prescribed-temperature wind; linear tides & Implemented; constant $\mu$ \\
H$\alpha$ transfer with supplied populations & Implemented \\
Multilevel ionization and internal radiation & Proposed; not implemented \\
Variable-$\mu$ feedback and energy balance & Deferred; not implemented \\
\bottomrule
\end{tabular}
\end{center}

Use cgs units: $r$ in cm, $v$ in cm\,s$^{-1}$, $\dot M$ in g\,s$^{-1}$,
$n_i,n_e,n_p,\nH$ in cm$^{-3}$, and $T$ in K. Public velocity APIs use
km\,s$^{-1}$; multiply by $10^5$ in the equations below. $\nH$ counts hydrogen
nuclei, $n_p$ protons, and $n_i$ neutral atoms in bound level $i$.
Neglect electron mass in $\rho=\mpart\nH$.
The mean molecular weight $\mu$ is in proton-mass units and is distinct from
the viewing-direction cosine $\mu_{\rm view}$.

Ordinary frequency is $\nu$ in Hz; the ExoJAX spectral grid is wavenumber
$\wn=\nu/c$ in cm$^{-1}$. Spectral densities transform as
$I_{\wn}=cI_\nu$ and $j_{\wn}=cj_\nu$, and normalized line profiles as
$\phi_{\wn}=c\phi_\nu$. Extinction at the same photon frequency is an inverse
length, not a spectral density. All radiation-rate integrals in
Section~\ref{sec:nlte} use ordinary frequency.

\section{Implemented population closure}\label{sec:populations}
\status{Implemented; prescribed ionization and excited reservoir}
The local functions are in \path{src/escape/hydrogen_populations.py}.
Callers must supply positive $r,v,T,T_\star,A_{ul},\wn_0,g_l,g_u$,
nonnegative $\dot M,n_e,\Upsilon,W$, and $x,f_{\rm exc}\in[0,1]$.
The population helper functions do not check or clip these physical bounds.
Steady spherical mass conservation gives
\begin{equation}
 \dot M=4\pi r^2\rho v,
 \qquad \nH=\frac{\dot M}{4\pi r^2v\mpart}.
 \tag{POP-1}\label{eq:continuity}
\end{equation}
For a lower level $l$ and upper level $u$, let
$\Delta E=hc\wn_0$, $g_l,g_u$ be statistical weights and $A_{ul}$ the
spontaneous rate in s$^{-1}$. The collision rates per atom are
\begin{equation}
 C_{ul}=n_e\frac{8.629\times10^{-6}\Upsilon_{ul}}{g_u\sqrt T},
 \qquad
 C_{lu}=\frac{g_u}{g_l}C_{ul}e^{-\Delta E/(\kb T)}.
 \tag{POP-2}\label{eq:collision-rates}
\end{equation}
Here $\Upsilon_{ul}$ is a dimensionless effective collision strength; the
de-excitation coefficient preceding $n_e$ has units cm$^3$\,s$^{-1}$ when
$T$ is in K. Maxwellian electrons and detailed balance are assumed.
For an externally prescribed mean intensity $J_\nu=W B_\nu(T_\star)$,
the dimensionless occupation is
\begin{equation}
 \eta=\frac{c^2J_{\nu_0}}{2h\nu_0^3}
     =\frac{W}{\exp[hc\wn_0/(\kb T_\star)]-1}.
 \tag{POP-3}\label{eq:occupation}
\end{equation}
$W\geq0$ is supplied; no attenuation or trapped radiation is solved.
Balancing the upward and downward rates gives
\begin{equation}
 n_l\left(C_{lu}+\frac{g_u}{g_l}A_{ul}\eta\right)
 =n_u\left[C_{ul}+A_{ul}(1+\eta)\right].
 \tag{POP-4}\label{eq:two-level-balance}
\end{equation}
Dividing and substituting Eq.~\eqref{eq:collision-rates} yields the implemented
\path{two_level_ratio}:
\begin{equation}
 R\equiv\frac{n_u}{n_l}
 =\frac{g_u}{g_l}
  \frac{C_{ul}e^{-\Delta E/(\kb T)}+A_{ul}\eta}
       {C_{ul}+A_{ul}(1+\eta)}.
 \tag{POP-5}\label{eq:two-level-ratio}
\end{equation}
For the current H$\alpha$ application, $l=2$, $u=3$, $g_l=8$ and $g_u=18$
represent unresolved shells. Model 3 supplies constant fractions
$x=n_p/\nH$ and $f_{\rm exc}$, the latter being the fraction of \emph{neutral}
hydrogen in the $n=2,3$ reservoir:
\begin{equation}
 n_e=x\nH,\qquad \mu=\frac{1}{1+x},\qquad
 n_2=\frac{f_{\rm exc}(1-x)\nH}{1+R},\qquad n_3=Rn_2.
 \tag{POP-6}\label{eq:reservoir}
\end{equation}
Thus the two-level ratio does not determine the ionized fraction or the
absolute excited population. The default $\Upsilon_{23}=35.52$ is frozen
at $kT=1$ eV; its table reconstruction and temperature dependence remain
open evidence tasks.

For collision domination, Eq.~\eqref{eq:two-level-ratio} approaches
$(g_u/g_l)e^{-\Delta E/(\kb T)}$. For $C_{ul}=0$, it approaches
$(g_u/g_l)\eta/(1+\eta)$. A Planck field at the gas temperature recovers the
Boltzmann ratio for any $C_{ul}$. These identities explain the saved tests;
they do not establish Saha ionization equilibrium or physical applicability.
The six-test snapshot uses $\Upsilon=10$, not the default 35.52.

\section{Implemented wind and temperature prescription}\label{sec:wind}
\status{Implemented; no energy equation}
Write $a^2=P/\rho=\kb T/(\mu\mpart)$. This is a pressure-derived thermal
speed, not the adiabatic sound speed with an additional $\gamma$ factor.
The momentum equation and mass conservation imply
\begin{align}
 v\frac{\dd v}{\dd r}
   &=-\frac{1}{\rho}\frac{\dd P}{\dd r}
     -\frac{GM_{\rm p}}{r^2}+\kappa r,
 &\kappa&=\frac{3GM_\star}{a_{\rm orb}^3},
 \tag{WIND-1}\label{eq:momentum}\\
 \left(v-\frac{a^2}{v}\right)\frac{\dd v}{\dd r}
   &=\frac{2a^2}{r}-\frac{\dd a^2}{\dd r}
     -\frac{GM_{\rm p}}{r^2}+\kappa r
     \equiv N(r).
 \tag{WIND-2}\label{eq:wind}
\end{align}
$\kappa$ has units s$^{-2}$ and defaults to zero. This is the one-dimensional
linear tide used in the target paper, with spherical continuity, not full
Roche geometry. For a prescribed $a^2(r)$, the selected regular transonic
branch has $v_c=a_c$ and $N(r_c)=0$. At constant $T,\mu$ and $\kappa=0$,
\begin{equation}
 r_c=\frac{GM_{\rm p}}{2a^2},\qquad
 \mathcal M^2-\ln\mathcal M^2
  =4\ln\frac{r}{r_c}+\frac{4r_c}{r}-3,
 \qquad \mathcal M=\frac{v}{a}.
 \tag{WIND-3}\label{eq:isothermal-wind}
\end{equation}
Choose the accelerating branch, subsonic below $r_c$ and supersonic above it.
The analytic isothermal solution initializes the numerical nonisothermal
solver; \path{nonisothermal.py} uses implicit roots for its converged wind
and sonic point.

The forward model prescribes
\begin{equation}
 T(r)=4500\,{\rm K}+(T_{\rm top}-4500\,{\rm K})
 \left[1-\exp\left(-\frac{r/\Rp-1}{0.12}\right)\right].
 \tag{WIND-4}\label{eq:temperature}
\end{equation}
In \texttt{nonisothermal} mode this temperature drives pressure, populations
and broadening. In \texttt{isothermal} mode the wind instead uses the supplied
$T_{\rm wind}$. Both modes currently hold $\mu$ constant. The spectral layers
cover $1$--$3\Rp$; the forward model's wind solve extends to $32\Rp$ and
requires a single regular sonic point within its domain. The initial
no-tide isothermal sonic point must also be inside that domain.

\section{Implemented spectral connection}\label{sec:spectrum}
\status{Implemented formal transfer with supplied populations}
For a profile normalized by $\int\phi_{\wn}\,\dd\wn=1$, the ExoJAX
\path{line_coefficients} convention is
\begin{align}
 \chi_{\wn}
 &=\frac{A_{ul}}{8\pi c\wn_0^2}
    \left(\frac{g_u}{g_l}n_l-n_u\right)\phi_{\wn},
 \tag{RT-1}\label{eq:extinction}\\
 j_{\wn}
 &=\frac{hc\wn_0A_{ul}}{4\pi}n_u\phi_{\wn},
 &S_{\wn}&=\frac{j_{\wn}}{\chi_{\wn}}
   =\frac{2hc^2\wn_0^3n_u}{(g_u/g_l)n_l-n_u}.
 \tag{RT-2}\label{eq:emissivity}
\end{align}
$\chi_{\wn}$ is in cm$^{-1}$; $j_{\wn}$ is in
erg\,s$^{-1}$\,cm$^{-3}$\,sr$^{-1}$ per cm$^{-1}$.
The prefactors use the rest line center (narrow-line approximation).
Absorption and emission share a profile. The source expression assumes
nonzero extinction; numerical zero-column limits are handled in the code.
The implemented isolated two-level closure has no population inversion for
finite positive temperature and nonnegative rates.

For a homogeneous layer, outward intensity satisfies
\begin{equation}
 I_{\wn}^{\rm out}=I_{\wn}^{\rm in}
       e^{-\Delta\tau_{\wn}/\mu_{\rm view}}
 +S_{\wn}\left(1-e^{-\Delta\tau_{\wn}/\mu_{\rm view}}\right),
 \qquad \Delta\tau_{\wn}=\chi_{\wn}\Delta r.
 \tag{RT-3}\label{eq:formal-transfer}
\end{equation}
The code integrates local plane-parallel columns over the visible dayside,
including projected radial, rotation and zonal velocities and instrumental
broadening. It uses a blackbody lower boundary, no atmospheric continuum
opacity, and fixed projected area at $\Rp$. Off-limb emission and the
observational reduction operator are omitted. These assumptions are separate
from the population model. Summed $n_2,n_3$ and shell-averaged line strengths
will need review when sublevels acquire non-statistical populations.

\section{Proposed multilevel NLTE problem}\label{sec:nlte}
\status{Proposed; not a closed or implemented numerical model}
The first stage takes $T(r)$, $\nH(r)$ and an external radiation field as
inputs, then replaces the external field by a coupled radiation calculation.
Resolve at least $1s,2s,2p$ and upper bound levels, plus protons and electrons.
The upper-level truncation and aggregation are not yet fixed. The reduced
hydrogen model of Garc\'ia Mu\~noz and Schneider~\cite{garcia} is a starting
reference, but its omission of bound--bound absorption other than Ly$\alpha$
cannot be carried over to the H$\alpha$ excitation problem here.

Let $P_{ij}$ be a bound--bound transition rate in s$^{-1}$, including
collisions and radiation. For a one-photon line, with Einstein $B$ defined
to multiply mean intensity rather than radiation energy density,
\begin{equation}
 \overline J_{ul}=\int\phi_{ul,\nu}J_\nu\,\dd\nu,\qquad
 P_{lu}=C_{lu}+B_{lu}\overline J_{ul},\qquad
 P_{ul}=C_{ul}+A_{ul}+B_{ul}\overline J_{ul}.
 \tag{SE-1}\label{eq:multilevel-rates}
\end{equation}
Here $C_{ij}$ can include electron and proton collisions. Two-photon decay
and sublevel mixing require their own rates; they are not represented by
inventing an ordinary $B$ coefficient for a forbidden line.

The net creation rate of bound level $i$ is
\begin{equation}
 Q_i=\sum_{j\ne i}(n_jP_{ji}-n_iP_{ij})
    +n_en_p\alpha_i^{\rm rad}
    +n_e^2n_pq_i^{(3)}
    -n_i\left(\Gamma_i+n_eq_i^{\rm ion}\right).
 \tag{SE-2}\label{eq:level-balance}
\end{equation}
$\alpha_i^{\rm rad}$ and $q_i^{\rm ion}$ are in cm$^3$\,s$^{-1}$;
$q_i^{(3)}$ is in cm$^6$\,s$^{-1}$, and $Q_i$ in cm$^{-3}$\,s$^{-1}$.
The radiative recombination coefficient may depend on $J_\nu$ through
stimulated recombination; inverse processes must be consistent before
claiming a full thermodynamic LTE limit. With threshold $\nu_i$ and
photoionization cross section $\sigma_i(\nu)$ in cm$^2$,
\begin{equation}
 \Gamma_i=4\pi\int_{\nu_i}^{\infty}
    \frac{\sigma_i(\nu)J_\nu}{h\nu}\,\dd\nu.
 \tag{SE-3}\label{eq:photoionization}
\end{equation}
This must include photoionization from excited states, particularly the
Balmer continuum. An EUV-only ground-state treatment is not the intended
KELT-9b model.

The initial local statistical-equilibrium approximation is
\begin{equation}
 Q_i=0\quad(i\ \hbox{bound}),\qquad
 \sum_i n_i+n_p=\nH,\qquad n_e=n_p.
 \tag{SE-4}\label{eq:conservation}
\end{equation}
Use the bound-level equations plus nuclear and charge conservation, or
replace one redundant species equation by nuclear conservation. Do not append
conservation to a redundant set of all bound and proton balance equations.
The solution must have nonnegative populations.

For steady advecting populations the fuller equation is
\begin{equation}
 \frac{1}{r^2}\frac{\dd(r^2vn_i)}{\dd r}
 =\nH v\frac{\dd(n_i/\nH)}{\dd r}=Q_i.
 \tag{SE-5}\label{eq:population-advection}
\end{equation}
Local equilibrium neglects this transport when relevant reaction times are
short compared with the flow time over population-gradient scales. That
ordering must be checked, not assumed from the existence of a steady wind.

For the coupled stage,
\begin{equation}
 \frac{\dd I_\nu}{\dd s}=-\chi_\nu I_\nu+j_\nu,
 \qquad J_\nu=\frac{1}{4\pi}\int I_\nu\,\dd\Omega.
 \tag{SE-6}\label{eq:radiation-feedback}
\end{equation}
The same populations determine the radiation coefficients and respond to
the resulting $J_\nu$. This internal mean intensity is additional to the
observer-directed formal spectrum. Ly$\alpha$ trapping, H$\alpha$ pumping,
continuum shielding and recombination radiation require explicit treatment
or a tested closure. An escape probability is a possible approximation,
not an established substitute for the target calculation.

\section{Residual and differentiation contract}\label{sec:implicit}
\status{Implemented for the prescribed wind; proposed for coupled NLTE}
After choosing grids, rates and boundary conditions, collect independent
unknowns into $y$ and model parameters into $\theta$. The actual discrete
equations, including conservation and radiation feedback, define
\begin{equation}
 F(y,\theta)=0,\qquad
 F_y\,\delta y=-F_\theta\,\delta\theta.
 \tag{AD-1}\label{eq:implicit-jvp}
\end{equation}
For a scalar objective $L$, a reverse-mode calculation may solve
\begin{equation}
 F_y^{\mathsf T}\lambda=\left(\frac{\partial L}{\partial y}\right)^{\mathsf T},
 \qquad
 \frac{\dd L}{\dd\theta}
 =\frac{\partial L}{\partial\theta}-\lambda^{\mathsf T}F_\theta.
 \tag{AD-2}\label{eq:implicit-vjp}
\end{equation}
These formulas require a sufficiently converged, locally isolated root and
an invertible reduced Jacobian. The solver and its derivative must use the
same residual, as required by JAX's \texttt{custom\_root}~\cite{jax}.
If radiation is eliminated by a transfer solve, its derivative still belongs
in $F_y$ and $F_\theta$. A local rate matrix evaluated at fixed $J_\nu$ is
not the derivative of the coupled model.
Use fixed discretizations initially, record residual and linear-solve errors,
and compare JVP/VJP results with independent finite differences before NUTS.

\section{Choices to record before implementation}\label{sec:choices}
The proposed system becomes a model only after specifying the level set,
atomic rates and their domains, upper-level cascade closure without double
counting, frequency redistribution, line broadening, geometry, irradiation,
and incoming radiation at both boundaries. A first benchmark should hold
$T$ and $\nH$ fixed and match these choices in the reference calculation.
Equation recording alone is not numerical or physical validation.

Later stages replace the prescribed $x$ by Eq.~\eqref{eq:conservation}, so
$\mu(r)=1/[1+n_p(r)/\nH(r)]$ varies. All derivatives of
$a^2=\kb T/(\mu\mpart)$ must then enter the wind residual, over the full wind
domain. A coupled energy/ionization model must also specify its own regularity
and boundary conditions; prescribed-profile wind conditions do not complete
that problem. H$^-$ continuum, heating/cooling and energy transport remain
deferred. The internal-energy convention must account consistently for
ionization and excitation energy; electron donors and metal cooling require
assessment before claiming a self-consistent thermal structure.

\section{Connection to MCEG}\label{sec:mceg}
\path{doc/model.tex} is the editable mathematical source, and
\path{doc/model.pdf} its rendered copy. Stable labels such as
\path{eq:two-level-ratio} identify equations even if page numbers change.
MCEG's \texttt{theory} Evidence links the files, label locations and hashes.
The formulation assessment covers only
Eqs.~\eqref{eq:continuity}--\eqref{eq:two-level-ratio} and their stated limits.
The wind, spectral and proposed NLTE sections are explanatory definitions,
not newly supported claims in that assessment.

The saved numerical assessment still refers only to the six population
tests at its original code revision. Physical accuracy remains unknown,
and inventory boundaries await human review. On a mathematical change,
review the model identity and assumptions, rebuild and inspect the PDF,
update the scoped formulation record, and obtain new numerical evidence
when the implemented equations change. Do not update old test evidence
merely because a document was added. MCEG does not parse or prove TeX equations;
it also does not build this PDF. Neither LaTeX nor MCEG is a runtime dependency
of \texttt{escape}.

\begin{thebibliography}{9}
\bibitem{zhang}
Y. Zhang et al. (2026), \emph{Hydrogen airglow from an escaping ultrahot
Jupiter atmosphere}, arXiv:2606.23899v1, Sections IV--V.
\url{https://arxiv.org/abs/2606.23899v1}.
\bibitem{garcia}
A. Garc\'ia Mu\~noz and P. C. Schneider (2019), \emph{Rapid escape of ultra-hot
exoplanet atmospheres driven by hydrogen Balmer absorption},
arXiv:1910.00267v1, Section II.
\url{https://arxiv.org/abs/1910.00267v1}.
\bibitem{jax}
JAX documentation, \texttt{jax.lax.custom\_root}, accessed 20 November 2026.
\url{https://docs.jax.dev/en/latest/_autosummary/jax.lax.custom_root.html}.
\bibitem{mceg}
MCEG v0.2 records, source revision \texttt{fff7afe}.
\url{https://github.com/HajimeKawahara/mceg/blob/fff7afe36f36627e002d27a953a53344ab22bad3/docs/en/v02.md}.
\end{thebibliography}
\end{document}
